\begin{figure}[t]
\centering
\begin{tikzpicture}[
  x=0.93cm,
  y=0.93cm,
  box/.style={draw, rounded corners=1.5pt, align=center, font=\scriptsize,
    inner sep=3.5pt, text width=2.15cm},
  native/.style={box, draw=green!45!black, very thick, fill=green!8},
  extended/.style={box, draw=blue!55!black, fill=blue!7},
  context/.style={box, draw=orange!70!black, fill=orange!9},
  design/.style={box, draw=violet!60!black, fill=violet!7, text width=2.55cm},
  flow/.style={-{Latex[length=1.8mm]}, semithick},
  nativeflow/.style={flow, draw=green!45!black, very thick},
  contextflow/.style={flow, dashed, draw=orange!70!black},
  note/.style={font=\tiny, align=left}
]

% Experimental design and contextual records.
\node[design] (design) at (1.55,2.05)
  {$2\!\times\!2$ design\\treatment: vehicle/active\\enrichment: standard/enriched\\four groups; two animals each};
\node[native] (animal) at (5.35,2.05)
  {animal-8\\\textbf{ISA Source}\\HCMO Subject};
\node[context] (housing) at (9.25,2.05)
  {cage + time-bounded\\HCMO housing\\assignment; 9 rows};
\draw[flow] (design) -- (animal);
\draw[contextflow] (housing) -- (animal);

% Native Source-to-Sample branch: kept left of the acquisition branch so that
% no arrow crosses another node or edge.
\node[native] (sample) at (0.95,0)
  {tissue specimen\\\textbf{ISA Sample}};
\node[native] (collection) at (3.25,0)
  {collection\\Protocol + Process};
\draw[nativeflow] (animal.south west) -- (collection.north east);
\draw[nativeflow] (collection) -- (sample);

% Full HCMO / extended ISA RO-Crate acquisition and analysis branch.
\node[extended] (recording) at (7.65,0)
  {recording\\Protocol + Process\\56 observations};
\node[extended] (raw) at (10.35,0)
  {raw ISA DataFile\\\texttt{dark-phase-}\\\texttt{activity.csv}};
\draw[flow] (animal.south east) -- (recording.north west);
\draw[flow] (recording) -- (raw);

\node[extended] (results) at (7.65,-2.05)
  {STATO-annotated\\result matrix +\\exact row fragments};
\node[extended] (model) at (10.35,-2.05)
  {linear mixed-model\\Protocol + Process};
\draw[flow] (raw) -- node[right, font=\tiny]{input} (model);
\draw[flow] (model) -- (results);

\node[note, text width=11.3cm, anchor=west] at (0,-3.45) {%
  \textcolor{green!45!black}{\rule{4mm}{0.7pt}} also generated by native ISA-API;
  \textcolor{blue!55!black}{\rule{4mm}{0.7pt}} full RDF/extended-crate path;
  \textcolor{orange!70!black}{\rule[0.5ex]{4mm}{0.5pt}} contextual link, not data production.\\[-1pt]
  \texttt{examples/isa-roundtrip/}: \texttt{README.md}; \texttt{canonical.ttl};
  \texttt{ro-crate-metadata.json};\\[-1pt]
  \texttt{data/dark-phase-activity.csv}; \texttt{data/model-results.csv};
  \texttt{native-isa/}; \texttt{loss/}.};

\end{tikzpicture}
\caption{\new{Executable 2 $\times$ 2 factorial fixture with repeated measures (seven daily observations per animal).} Each entity and output has a distinct role. The green
animal Source $\rightarrow$ collection Process $\rightarrow$ genuine tissue
Sample path is reproducibly generated as ISA-JSON and ISA-Tab. The remaining
factor, housing, whole-animal recording, repeated-observation, and semantic
STATO/file-fragment structures are validated in HCMO RDF and the extended
ISA RO-Crate, but remain declared losses in the native ISA projection. The
result CSV exposes the reviewed STATO IRIs and labels while the RDF retains the
typed entities and exact row-fragment links.}
\label{fig:roundtrip}
\end{figure}
